% =====================================================================
% Rozdzial 1: Wprowadzenie
% =====================================================================

\chapter{Wprowadzenie}
\label{chap:wprowadzenie}

\section{Cel modelu LEM}
\label{sec:cel-lem}
Informacja o lokalizacji jest wykorzystywana przez wiele systemów technicznych działających w różnych domenach zastosowań. Systemy zarządzania infrastrukturą, automatyzacji, bezpieczeństwa, dokumentacji technicznej oraz usług lokalizacyjnych wymagają jednoznacznego opisu położenia obiektów fizycznych. W praktyce informacje te są często reprezentowane przy użyciu lokalnych konwencji zależnych od systemu, producenta lub technologii. Stosowane podejścia obejmują między innymi:

\begin{itemize}
    \item nieustrukturyzowane opisy tekstowe w systemach CMDB, bazach danych lub dokumentacji technicznej,
    \item lokalne schematy nazewnictwa urządzeń i przestrzeni,
    \item oznaczenia wynikające z dokumentacji architektonicznej lub organizacyjnej,
    \item informacje utrzymywane jako wiedza operacyjna określonych osób lub zespołów.
\end{itemize}

Takie sposoby opisu często nie zapewniają jednoznacznej, deterministycznej i maszynowo interpretowalnej informacji. Konsekwencją braku wspólnego modelu semantycznego są m.in. ograniczona interoperacyjność, utrudniona automatyzacja oraz problemy z utrzymaniem spójności.

Model Kodowania Lokalizacji (\textit{Location Encoding Model}, LEM) definiuje wspólny model semantyczny, którego celem jest zapewnienie jednoznacznego znaczenia komponentów niezależnie od sposobu ich reprezentacji czy implementacji. LEM definiuje znaczenie informacji, podczas gdy sposób jej wyrażenia w konkretnym formacie pozostaje warstwą niezależną.

\section{Zakres dokumentu}
\label{sec:zakres}
Niniejszy dokument definiuje \textit{Location Encoding Model} (LEM) jako model informacji przeznaczony do opisu lokalizacji. Specyfikacja określa:

\begin{itemize}
    \item model semantyczny informacji o lokalizacji,
    \item komponenty modelu LEM oraz ich znaczenie,
    \item zasady interpretacji informacji, równoważności semantycznej oraz spójności modelu,
    \item zasady reprezentacji informacji LEM oraz tworzenia profili zastosowań.
\end{itemize}

Dokument nie definiuje formatów wymiany danych, protokołów komunikacyjnych, technologii pomiaru lokalizacji, algorytmów pozycjonowania ani wymagań specyficznych dla pojedynczego środowiska.

LEM nie definiuje mechanizmów identyfikacji obiektów, pozyskiwania informacji lokalizacyjnej ani metod określania położenia. Model opisuje wyłącznie semantykę informacji o lokalizacji już dostępnej.

\section{Wkład dokumentu}
\label{sec:wklad}
Wkład LEM obejmuje cztery elementy:
\begin{enumerate}
    \item niewielki rdzeń semantyczny oddzielający znaczenie lokalizacji od jej reprezentacji;
    \item deterministyczną postać kanoniczną i jawne reguły równoważności opisów częściowych;
    \item spójne odwzorowania na JSON, TLV i YANG oraz mapowania integracyjne do Civic Address i GeoJSON;
    \item mechanizm profili umożliwiający specjalizację domenową bez redefiniowania rdzenia modelu.
\end{enumerate}

\section{Niezależność semantyki od reprezentacji}
\label{sec:niezaleznosc}
Model LEM rozdziela znaczenie informacji od sposobu jej wyrażenia. Znaczenie jest określone przez model semantyczny i interpretację kanoniczną. Ta sama informacja może być wyrażona w różnych reprezentacjach (tekstowej, strukturalnej, binarnej), które są semantycznie równoważne, jeśli prowadzą do tej samej interpretacji zgodnie z zasadami LEM.

\section{Struktura dokumentu}
\label{sec:struktura-dok}
\begin{description}
    \item[Rozdział 2:] Definiuje model semantyczny LEM, komponenty oraz zasady interpretacji.
    \item[Rozdział 3:] Definiuje zasady reprezentacji modelu w różnych formatach danych.
    \item[Rozdział 4:] Definiuje profile zastosowań dla specjalizacji domenowych.
    \item[Rozdziały 5--7:] Omawiają implementację, bezpieczeństwo, prywatność i ewolucję modelu.
    \item[Rozdział 8:] Osadza LEM względem istniejących standardów i modeli.
\end{description}

\section{Terminologia i konwencje}
\label{sec:terminologia}
Na potrzeby dokumentu przyjęto następujące definicje:

\begin{description}
    \item[Model informacji (\textit{Information Model}):] Abstrakcyjny model definiujący strukturę, znaczenie oraz relacje elementów.
    \item[Informacja o lokalizacji:] (\textit{Location Information}) Informacja opisana przy użyciu komponentów LEM.
    \item[Komponent modelu (\textit{Model Component}):] Element LEM posiadający znaczenie semantyczne.
    \item[Przestrzeń referencyjna (\textit{Reference Space}):] Przestrzeń odniesienia dla interpretacji przestrzennej.
    \item[Położenie względne (\textit{Relative Position}):] Relacja przestrzenna względem przestrzeni referencyjnej.
    \item[Reprezentacja (\textit{Representation}):] Sposób wyrażenia informacji przy użyciu formatu danych.
    \item[Interpretacja kanoniczna (\textit{Canonical Interpretation}):] Jednoznaczne znaczenie wynikające z modelu LEM.
    \item[Profil LEM (\textit{LEM Profile}):] Specjalizacja modelu określająca wymagania dla danego zastosowania.
\end{description}

Słowa kluczowe zapisane wielkimi literami są używane zgodnie z BCP~14 \cite{RFC2119,RFC8174}: \textbf{MUSI} odpowiada \textit{MUST}, \textbf{NIE MOŻE} --- \textit{MUST NOT}, \textbf{POWINIEN} --- \textit{SHOULD}, a \textbf{MOŻE} --- \textit{MAY}. Przykłady mają charakter ilustracyjny, o ile nie oznaczono ich inaczej.
